% ============================================================================
% Fichier  : partie4/P04_C13_audit_gouvernance_dsi.tex
% Titre    : Auditer la Gouvernance et la Fonction Informatique
% Auteur   : MINKA MI NGUIDJOÏ Thierry Emmanuel
% Version  : 1.0 -- novembre 2026
% Base   : ACSI C13, repris et développé pour INF3226
% ============================================================================
% Sources : guide d'audit des SI du CHAI (2014), fiche 3.1 « Audit du
% pilotage des systèmes d'information » (six axes : rôle et positionnement,
% planification stratégique, budgets et coûts, mesure de la performance,
% organisation de la DSI, cadre légal ; distinction contenant / contenu),
% paraphrasée ; COBIT 2019 (EDM01, APO01, APO02, APO05, APO06, APO07, APO09,
% MEA01) comme source de critères ; décomposition du domaine A de LiZbiZ
% cohérente avec le chapitre 9 (16 / 3 / 1). La gouvernance elle-même est
% traitée dans les UE MSN-2S3-MSI01 et MSN-2S3-MSI03 (GCI).
% ============================================================================

\chapter{Auditer la Gouvernance et la Fonction Informatique}
\label{chap:audit-gouvernance}
\chaptermark{Auditer la gouvernance}

\epigraph{%
    « L'auditeur ne juge pas la stratégie.\\
    Il juge la manière dont on l'a décidée, suivie et corrigée. »%
}{Règle de l'audit de gouvernance}

\niveauChapitre{Intermédiaire}{%
    Maîtrise du Chapitre~\ref{chap:conclure} (mesures et observation) et du
    Chapitre~\ref{chap:referentiels} (critères tirés de COBIT)%
}

\begin{boxobjectifs}
\begin{itemize}
    \item Décomposer l'audit de la fonction informatique en sous-domaines~:
          rôle et organes, planification stratégique, budgets, performance,
          organisation et compétences.
    \item Distinguer, pour chaque document de gouvernance, le contrôle du
          \emph{contenant} (processus d'élaboration et de validation) et
          celui du \emph{contenu} (pertinence, cohérence).
    \item Évaluer les organes de gouvernance des SI à partir de la trace de
          leurs décisions.
    \item Auditer un indicateur de pilotage~: définition, source, calcul,
          reproduction.
\end{itemize}
\end{boxobjectifs}

\bigskip

% ============================================================================
\section*{Question directrice : l'informatique de LiZbiZ est-elle bien gouvernée~?}
% ============================================================================

Au Chapitre~\ref{chap:conclure}, le domaine~A de LiZbiZ, la gouvernance du
SI, sortait en conformité claire~: 16~\OUI{}, 3~\NON{}, 1~\NEK{}. Le conseil
d'administration s'en félicite et demande à AUDIX CONSEIL s'il peut
« passer à autre chose ». M.~BEYALA répond par une question~: qu'y a-t-il
derrière ces vingt réponses~?

\medskip
\textbf{Plan du chapitre.}
\begin{enumerate}
    \item Décomposer l'audit de la fonction informatique (section~\ref{sec:c17-decomposer}).
    \item Le rôle, le positionnement et les organes de gouvernance (section~\ref{sec:c17-organes}).
    \item La planification stratégique et les budgets (section~\ref{sec:c17-planification}).
    \item L'organisation et les compétences (section~\ref{sec:c17-organisation}).
    \item Auditer un indicateur de pilotage (section~\ref{sec:c17-indicateurs}).
    \item Auditer un engagement de service (section~\ref{sec:c17-sla}).
\end{enumerate}

\separateur

% ============================================================================
\section{Décomposer l'Audit de la Fonction Informatique}
\label{sec:c17-decomposer}
% ============================================================================

\subsection{Ce que l'on audite, et ce que l'on n'audite pas}

L'audit de la gouvernance porte sur des décisions~: qui décide, selon
quelle procédure, avec quelle information, et comment la décision est
suivie. L'auditeur n'a pas à dire si LiZbiZ aurait dû choisir tel
progiciel ou ouvrir tel centre de données~: ce choix appartient aux
dirigeants. Il vérifie que le choix a été préparé, décidé par l'organe
compétent, tracé, et suivi.

Le guide du CHAI formule cette distinction pour les documents de
pilotage~: on analyse d'abord le \emph{contenant}, c'est-à-dire le
processus d'élaboration et de validation, puis le \emph{contenu}, sa
pertinence et son alignement~\cite{chai2014guide}.

\begin{boxdef}
\termecle{Contenant et contenu}\index{contenant et contenu}~: pour un
document de gouvernance (schéma directeur, budget, tableau de bord), le
\emph{contenant} est la manière dont il a été élaboré, validé et mis à
jour~; le \emph{contenu} est ce qu'il dit, et sa cohérence avec les
objectifs de l'organisation. Le contenant se teste par des preuves
(procès-verbaux, versions, signatures)~; le contenu se juge par
rapprochement avec d'autres documents.
\end{boxdef}

\subsection{Cinq sous-domaines}

En reprenant les axes du guide du CHAI et les objectifs de COBIT~2019,
AUDIX CONSEIL décompose le domaine~A en cinq sous-domaines de quatre
contrôles chacun. Le cadre légal, sixième axe du CHAI, est traité dans les
domaines où chaque obligation s'applique (Chapitre~\ref{chap:droit}).

\begin{tableaucours}{cL{3.4cm}L{5.4cm}L{2.4cm}}{%
    \tableauHeader{Code} &
    \tableauHeader{Sous-domaine} &
    \tableauHeader{Contrôles} &
    \tableauHeader{Critères COBIT}
}
A.1 & Rôle, positionnement et organes & Rattachement de la DSI à la
direction générale~; mission écrite~; comité de pilotage du SI actif~;
rôle des métiers dans les projets & EDM01, APO01 \\
A.2 & Planification stratégique & Schéma directeur approuvé~; revue
annuelle~; projets rattachés au schéma~; ressources adaptées & APO02,
APO05 \\
A.3 & Budgets et coûts & Processus budgétaire écrit~; budget exhaustif~;
suivi des écarts~; bilan économique des projets & APO06 \\
A.4 & Performance et indicateurs & Objectifs assignés à la DSI~; tableau
de bord~; engagements de service~; satisfaction des utilisateurs &
MEA01, APO09 \\
A.5 & Organisation et compétences & Organigramme et fiches de poste~;
séparation études et exploitation~; dépendance aux personnes~; plan de
formation & APO07 \\
\end{tableaucours}
\vspace{-0.8em}
\begin{center}{\small\itshape Tableau~17.1. Décomposition du domaine~A,
gouvernance et fonction informatique}\end{center}

\begin{boxlecture}{Domaine A décomposé}
Les vingt contrôles du domaine~A, lus par sous-domaine~:

\medskip
\begin{tableaucours}{L{4.2cm}cccccc}{%
    \tableauHeader{Sous-domaine} &
    \tableauHeader{\OUI{}} & \tableauHeader{\NON{}} & \tableauHeader{\NEK{}} &
    \tableauHeader{$\mu$} & \tableauHeader{$U$} & \tableauHeader{$H$}
}
A.1 Rôle et organes & 4 & 0 & 0 & 1,00 & 0,00 & 0,00 \\
A.2 Planification & 3 & 1 & 0 & 0,75 & 0,00 & 0,81 \\
A.3 Budgets et coûts & 3 & 0 & 1 & 0,75 & 0,25 & 0,81 \\
A.4 Performance & 3 & 1 & 0 & 0,75 & 0,00 & 0,81 \\
A.5 Organisation & 3 & 1 & 0 & 0,75 & 0,00 & 0,81 \\
Domaine A & 16 & 3 & 1 & 0,80 & 0,05 & 0,88 \\
\end{tableaucours}

\medskip
\textbf{Constatation.} Les trois \NON{} portent sur la revue annuelle du
schéma directeur (A.2), les engagements de service (A.4) et la dépendance
à une seule personne (A.5). Le \NEK{} porte sur le bilan économique des
projets (A.3).

\medskip
\textbf{Analyse} (confiance moyenne). Le domaine est en conformité claire,
mais le sous-domaine A.3 est opaque~: $U = 0{,}25$. Agréger le domaine
dilue une incertitude locale, comme au Chapitre~\ref{chap:decomposition}.
Et le \NON{} de A.5 rejoint un constat d'un autre domaine~: le responsable
informatique, M.~KOUNOU, détient tous les droits de l'application
(Chapitre~\ref{chap:controle-interne}). Une bonne gouvernance sur le
papier n'empêche pas la concentration du pouvoir sur une personne.

\medskip
\hfill\lecture{0{,}80}{0{,}05}{0{,}88}
\end{boxlecture}

% ============================================================================
\section{Rôle, Positionnement et Organes de Gouvernance}
\label{sec:c17-organes}
% ============================================================================

\subsection{Le positionnement de la DSI}

Le guide du CHAI ouvre sa fiche de pilotage par le positionnement de
l'informatique~: la DSI est-elle rattachée à la direction générale et
présente au comité de direction~? Un document définit-il son rôle et son
périmètre~? Les métiers jouent-ils leur rôle de maîtrise d'ouvrage
(besoins, priorités, recette), ou laissent-ils la DSI décider à leur
place~?

Un signal revient souvent dans les organisations camerounaises~: la DSI
pilote seule des projets métier, parce que les directions métier ne
désignent pas de responsable. Le risque n'est pas technique~: c'est un
système conforme au cahier des charges et inutile aux utilisateurs.

\subsection{Évaluer un organe par la trace de ses décisions}

Un comité de pilotage du SI, un comité des risques, un comité d'audit
existent souvent sur le papier. L'auditeur ne s'arrête pas à leur
existence~: il évalue ce qu'ils décident.

\begin{boxprocedure}
\textbf{Évaluer un organe de gouvernance des SI.}
\begin{enumerate}
    \item \textbf{Mandat}~: un texte fixe sa composition, ses attributions et
          sa fréquence de réunion. (Preuve~: charte du comité, décision
          de création.)
    \item \textbf{Fonctionnement}~: il s'est réuni à la fréquence prévue,
          avec le quorum. (Preuve~: procès-verbaux de la période.)
    \item \textbf{Information}~: il a reçu les éléments utiles pour décider~:
          risques, avancement des projets, incidents majeurs. (Preuve~:
          dossiers de séance.)
    \item \textbf{Décision}~: ses procès-verbaux contiennent des décisions,
          avec un responsable et une échéance, et non de simples «~points
          d'information~».
    \item \textbf{Suivi}~: les décisions sont reprises à la séance
          suivante, et leur exécution est vérifiable.
\end{enumerate}
\end{boxprocedure}

\begin{boxscenario}{Opération SANAGA. Le comité des risques et la bascule}
\textbf{Faits.} Mme~EKAMBI lit les procès-verbaux du comité des risques de
la BANQUE SANAGA pour 2025 et le premier trimestre 2026. Le comité s'est
réuni quatre fois, conformément à sa charte. Aucun dossier de séance ne
mentionne la bascule du système central, ni le risque de panne coté
« 3 sur 5 » par la DSI (Chapitre~\ref{chap:risque}).

\medskip
\textbf{Cotation.} Mandat~: \OUI{}. Fonctionnement~: \OUI{}.
Information~: \NON{}. Décision et suivi sur le risque informatique~: sans
objet faute d'information, donc \NON{} au regard du critère « le comité
statue sur les risques majeurs, y compris informatiques ».

\medskip
\textbf{Analyse} (confiance élevée). L'organe existe et fonctionne, mais
il n'a jamais été saisi du risque qui s'est réalisé. La défaillance n'est
pas dans le comité~: elle est dans le circuit qui aurait dû l'informer.
La recommandation porte sur l'échelle commune de cotation et sur la
remontée des risques informatiques, pas sur le comité lui-même.
\end{boxscenario}

% ============================================================================
\section{Planification Stratégique et Budgets}
\label{sec:c17-planification}
% ============================================================================

\subsection{Le schéma directeur}

Le schéma directeur, ou plan stratégique du SI, est le premier document de
strate~1 (Chapitre~\ref{chap:planifier}). Il s'audite en deux temps.

\begin{tableaucours}{L{2.4cm}L{5.0cm}L{4.8cm}}{%
    \tableauHeader{Angle} &
    \tableauHeader{Questions} &
    \tableauHeader{Preuves}
}
Contenant & Qui l'a élaboré~? Qui l'a approuvé~? Est-il revu au moins une
fois par an~? & Versions datées, procès-verbal d'approbation, comptes
rendus de revue \\
Contenu & Les projets en cours y figurent-ils~? Les priorités
correspondent-elles à la stratégie de l'organisation~? Les ressources
suffisent-elles~? & Rapprochement avec le portefeuille de projets, le
plan stratégique de l'entreprise, le budget \\
\end{tableaucours}
\vspace{-0.8em}
\begin{center}{\small\itshape Tableau~17.2. Auditer un schéma directeur}\end{center}

Un test simple donne beaucoup d'information~: prendre la liste des projets
informatiques en cours et chercher chacun dans le schéma directeur. Un
projet absent du schéma a été décidé hors du circuit~; un projet du
schéma qui n'a jamais démarré signale un plan déconnecté des moyens.

\begin{boxinfo}
\textbf{Le contrôle A.2.2 de LiZbiZ.} Le schéma directeur 2025-2028 est
approuvé par le conseil (\OUI{} pour A.2.1). Mais aucune revue n'a eu
lieu depuis son approbation, alors que l'ouverture du centre de données de
Bafoussam en a modifié les priorités~: \NON{} pour A.2.2, revue annuelle.
\end{boxinfo}

\subsection{Les budgets et les coûts}

L'auditeur examine le processus budgétaire (qui propose, qui arbitre, qui
approuve), le périmètre du budget (tous les sites, charges et
investissements, dépenses informatiques engagées hors DSI), et le suivi
des écarts. Le guide du CHAI recommande de manier les ratios de coûts et
les comparaisons avec d'autres organisations avec précaution~: un ratio
n'est fiable que si les périmètres sont comparables~\cite{chai2014guide}.

Le bilan économique d'un projet compare les coûts et les bénéfices
constatés à ceux de l'étude préalable. Chez LiZbiZ, le DAF affirme qu'un
bilan de l'application intégrée a été présenté au conseil en 2025, mais
ne le retrouve pas~: \NEK{} pour A.3.4, en attendant le procès-verbal.

% ============================================================================
\section{Organisation et Compétences}
\label{sec:c17-organisation}
% ============================================================================

L'organisation de la DSI s'audite avec les outils du
Chapitre~\ref{chap:controle-interne}~: organigramme, fiches de poste,
séparation des tâches. Trois points méritent l'attention de l'auditeur
des SI.

\begin{itemize}
    \item \textbf{La séparation entre études et exploitation}~: ceux qui
          développent ne doivent pas accéder aux données ni aux programmes
          de production. Le Chapitre~\ref{chap:audit-continuite} y revient.
    \item \textbf{La dépendance aux personnes}~: une application connue
          d'une seule personne, un compte d'administration détenu par une
          seule personne, sont des risques de continuité et de fraude. Le
          guide du CHAI invite à apprécier cette dépendance au regard de
          l'âge des applications, de la taille des équipes et de la
          documentation.
    \item \textbf{Les compétences}~: un plan de formation nominatif, et un
          effort de formation suivi sur plusieurs années.
\end{itemize}

\begin{boxalert}
\textbf{Violation typique.} Conclure que la fonction informatique est bien
organisée parce que l'organigramme est à jour et les fiches de poste
signées.

\smallskip
\textbf{Règle~:} on confronte l'organigramme à la matrice des droits. Chez
LiZbiZ, l'organigramme sépare l'administration des droits et la gestion~;
la matrice donne tout à M.~KOUNOU. C'est la matrice qui dit ce qui se
passe (Chapitre~\ref{chap:planifier}, strates documentaires).
\end{boxalert}

% ============================================================================
\section{Auditer un Indicateur de Pilotage}
\label{sec:c17-indicateurs}
% ============================================================================

\subsection{Un indicateur est une affirmation}

Un tableau de bord présente des chiffres~: taux de disponibilité, délai de
résolution des incidents, avancement des projets. Pour l'auditeur, chacun
est une \emph{affirmation} de la direction, à vérifier comme toute autre.
Un indicateur faux est pire que l'absence d'indicateur~: il rassure à tort
le comité qui le lit.

\begin{boxprocedure}
\textbf{Auditer un indicateur.}
\begin{enumerate}
    \item \textbf{Définition}~: une fiche décrit l'indicateur, sa formule,
          sa source, sa fréquence, son responsable et sa cible.
    \item \textbf{Source}~: les données d'entrée viennent d'un système, et
          non d'une saisie manuelle invérifiable.
    \item \textbf{Reproduction}~: l'auditeur extrait les données sources
          (Chapitre~\ref{chap:preuve}) et recalcule l'indicateur selon la
          formule.
    \item \textbf{Écart}~: il explique tout écart entre la valeur publiée et
          la valeur recalculée.
    \item \textbf{Usage}~: l'indicateur est lu par l'organe compétent, et un
          dépassement de cible déclenche une action tracée.
\end{enumerate}
\end{boxprocedure}

\begin{boxlecture}{Disponibilité affichée}
Le tableau de bord de la DSI de LiZbiZ affiche, pour avril~2026, une
disponibilité de l'application intégrée de 99,6~\%, au-dessus de la cible
de 99,5~\%. Avril compte 720~heures. AUDIX CONSEIL extrait le journal des
incidents~: 7,9~heures d'indisponibilité, dont 5,0~heures classées
« maintenance planifiée » et exclues du calcul. Aucune de ces
maintenances n'a été annoncée aux utilisateurs.

\medskip
\textbf{Calcul.} Valeur publiée~: $(720 - 2{,}9) / 720 \approx 99{,}6~\%$.
Valeur recalculée selon la définition (seules les maintenances annoncées
sont exclues)~: $(720 - 7{,}9) / 720 \approx 98{,}9~\%$, sous la cible.

\medskip
\textbf{Analyse} (confiance élevée). L'indicateur n'est pas faux par erreur
de calcul~: il est faussé par un classement. Le contrôle A.4.2
(« le tableau de bord est fiable ») passe de \OUI{} à \NON{}. Le
domaine~A devient 15~\OUI{}, 4~\NON{}, 1~\NEK{}~: il reste en conformité
claire, mais l'écart se concentre désormais dans le pilotage.

\medskip
\hfill\lecture{0{,}75}{0{,}05}{0{,}99}
\end{boxlecture}

\begin{boxmetier}{Auditeur SI}
\textbf{Recalculer un indicateur}
\begin{enumerate}
    \item Demander la fiche de l'indicateur~; à défaut, faire écrire la
          définition par son responsable avant tout calcul.
    \item Extraire les données sources avec une lettre d'extraction et une
          empreinte (Chapitre~\ref{chap:preuve}).
    \item Recalculer, dans un classeur conservé au dossier, chaque étape de
          la formule.
    \item Lister les exclusions et vérifier que chacune est prévue par la
          définition.
    \item Présenter l'écart au responsable avant de conclure~: c'est le
          contradictoire (Chapitre~\ref{chap:mission}).
\end{enumerate}
\end{boxmetier}

\begin{boxmetier}{Gestionnaire de risques}
\textbf{Faire remonter le risque informatique}
\begin{enumerate}
    \item Coter les risques informatiques sur l'échelle commune de
          l'entreprise (Chapitre~\ref{chap:risque}).
    \item Inscrire au dossier du comité des risques tout risque informatique
          au-dessus du seuil d'appétence, et tout projet qui modifie un
          système critique.
    \item Vérifier, à chaque séance, que les décisions sur ces risques ont
          un responsable et une échéance.
    \item Rapprocher chaque trimestre le registre de la DSI et le registre
          de l'entreprise~: un risque présent dans le premier et absent du
          second est un écart.
\end{enumerate}
\end{boxmetier}

% ============================================================================

% ============================================================================
\section{Auditer un Engagement de Service}
\label{sec:c17-sla}
% ============================================================================

\subsection{Ce qu'est un engagement de service}

\begin{boxdef}
\termecle{Engagement de service}\index{engagement de service}
(\textit{service level agreement}, SLA)~: accord écrit entre un
fournisseur de service (la DSI, ou un prestataire) et ses utilisateurs,
qui fixe des niveaux mesurables (disponibilité, délai de résolution), la
manière de les mesurer et les conséquences d'un manquement.
\end{boxdef}

Un engagement de service se décompose comme une fiche d'assertion~: un
indicateur (critère), une méthode de calcul (test), une source de données
(trace). L'auditeur vérifie qu'il existe, qu'il est mesuré comme il est
écrit, et qu'il est respecté.

\subsection{Recalculer la disponibilité}

La disponibilité d'un service sur une période se calcule ainsi~:
\[
    \text{disponibilité} = \frac{T_{\text{ouverture}} - T_{\text{indisponibilité}}}{T_{\text{ouverture}}},
\]
où $T_{\text{ouverture}}$ est la durée pendant laquelle le service devait
être disponible. Tout l'enjeu de l'audit tient dans les exclusions~: les
« maintenances programmées » sont-elles déduites de l'indisponibilité, et
étaient-elles réellement programmées~?

\begin{boxprocedure}
\textbf{Auditer un indicateur de disponibilité}
\begin{enumerate}
    \item Obtenir la définition écrite de l'indicateur (engagement de
          service, tableau de bord) et la période d'ouverture.
    \item Extraire la liste des incidents et des maintenances de la période,
          avec leurs heures de début et de fin.
    \item Vérifier que chaque maintenance déduite a été annoncée avant son
          début (trace~: avis aux utilisateurs, ticket).
    \item Recalculer la disponibilité, avec et sans les exclusions
          contestées.
    \item Comparer au chiffre publié, et expliquer l'écart.
\end{enumerate}
\end{boxprocedure}

\begin{boxinfo}
\textbf{Rappel chiffré.} La section~\ref{sec:c17-indicateurs} a montré, pour LiZbiZ, une
disponibilité publiée de 99,6~\% en avril, recalculée à 98,9~\% après
réintégration de 5,0~heures de « maintenance » non annoncée. Sur 720~heures
d'ouverture, un écart de 0,7~point représente environ 5~heures~: c'est la
taille d'un incident majeur rendu invisible par une exclusion.
\end{boxinfo}

\separateur

\begin{boxresume}
\textbf{Ce qu'il faut retenir de ce chapitre~:}
\begin{itemize}
    \item \textbf{L'audit de la gouvernance} porte sur les décisions~: qui
          décide, comment, avec quelle information, avec quel suivi~; il ne
          juge pas l'opportunité des choix stratégiques.
    \item \textbf{Cinq sous-domaines}~: rôle et organes, planification,
          budgets, performance, organisation~; le cadre légal se traite
          dans chaque domaine concerné.
    \item \textbf{Contenant, puis contenu}~: la manière dont un document a
          été élaboré et approuvé, puis sa pertinence.
    \item \textbf{Un organe s'évalue} par la trace de ses décisions~:
          mandat, fonctionnement, information, décision, suivi.
    \item \textbf{Un indicateur est une affirmation}~: on le recalcule à
          partir des données sources, et on explique chaque écart.
    \item \textbf{Agréger un domaine} peut masquer une opacité locale~:
          lire aussi les sous-domaines.
    \item \textbf{Un engagement de service} s'audite comme une fiche d'assertion~: définition, mesure, exclusions annoncées, recalcul.
\end{itemize}
\end{boxresume}

\bigskip

\begin{boxexo}[title={\ding{45}\quad Exercice~17.1. Classer les points de contrôle \niveaubase}]
Rattachez chaque point de contrôle à un sous-domaine (A.1 à A.5) et dites
s'il porte sur le contenant ou sur le contenu.
\begin{enumerate}[(a)]
    \item Le schéma directeur a été approuvé par le conseil en 2025.
    \item Le budget informatique inclut les dépenses informatiques des
          agences.
    \item Le comité de pilotage du SI s'est réuni quatre fois en 2025.
    \item Les projets en cours figurent tous au schéma directeur.
    \item Chaque informaticien a un plan de formation nominatif.
\end{enumerate}
\end{boxexo}

\begin{boxexo}[title={\ding{45}\quad Exercice~17.2. Recalculer un délai de résolution \niveaumoyen}]
Le tableau de bord de la DSI de LiZbiZ annonce que 92~\% des incidents de
mai~2026 ont été résolus en moins de 8~heures. Le journal compte
150~incidents~: 126 résolus en moins de 8~heures, 11 résolus en plus de
8~heures, et 13 incidents « reclassés en demande de service » après
24~heures, exclus du calcul.
\begin{enumerate}[(a)]
    \item Retrouvez la valeur publiée.
    \item Recalculez l'indicateur en réintégrant les incidents reclassés.
    \item Que faut-il demander pour savoir si le reclassement est justifié~?
          Quelle cotation retenir en attendant~?
    \item Rédigez la constatation et l'analyse de l'observation.
\end{enumerate}
\end{boxexo}

\begin{boxexo}[title={\ding{45}\quad Exercice~17.3. Le comité des risques de la BANQUE SANAGA \niveauexpert}]
Reprenez le scénario « Le comité des risques et la bascule ».
\begin{enumerate}[(a)]
    \item Rédigez l'observation en sept éléments (Chapitre~\ref{chap:conclure}).
    \item Pourquoi la recommandation ne vise-t-elle pas le comité
          lui-même~?
    \item Quels contrôles ajouteriez-vous au sous-domaine A.1 de la banque
          pour couvrir ce défaut à l'avenir~? Rédigez la fiche d'assertion
          de l'un d'eux.
    \item Quel indicateur proposeriez-vous au comité d'audit pour suivre la
          remontée des risques informatiques~?
\end{enumerate}
\end{boxexo}

\bigskip

\begin{boxexo}[title={\ding{45}\quad Exercice~17.4. Les organes \niveaubase}]
Pour un comité de pilotage informatique, citez les cinq points
d'évaluation d'un organe (mandat, fonctionnement, information, décision,
suivi) et la pièce qui permet de vérifier chacun.
\end{boxexo}

\begin{boxexo}[title={\ding{45}\quad Exercice~17.5. Un schéma directeur \niveaumoyen}]
Le schéma directeur de LiZbiZ (2024-2027) n'a jamais été revu. Quelle
fiche d'assertion rédigez-vous~? Quelle cote, et quelle conséquence pour
le sous-domaine A.2~?
\end{boxexo}

\begin{boxexo}[title={\ding{45}\quad Exercice~17.6. La dépendance à une personne \niveaumoyen}]
M.~KOUNOU est le seul à connaître l'administration de la base de
l'application intégrée. Rédigez la constatation, le risque et deux
recommandations, en évitant toute mise en cause personnelle.
\end{boxexo}

\begin{boxexo}[title={\ding{45}\quad Exercice~17.7. Recalculer une disponibilité \niveaumoyen}]
En mai~2026, le système de facturation de LiZbiZ devait être ouvert
22~jours de 7~h à 19~h. Le tableau de bord publie 99,8~\%. Les journaux
montrent 3~h d'incidents et 4~h de « maintenance », dont 2~h annoncées.
Recalculez la disponibilité selon l'engagement (seules les maintenances
annoncées sont exclues) et concluez.
\end{boxexo}

\begin{boxexo}[title={\ding{45}\quad Exercice~17.8. L'engagement de DIGICORE \niveaumoyen}]
Le contrat de DIGICORE AFRIQUE prévoit une disponibilité de 99,5~\% par
mois, sans dire comment elle est mesurée ni par qui. Quel est le défaut, et
quelle clause proposez-vous~?
\end{boxexo}

\begin{boxexo}[title={\ding{45}\quad Exercice~17.9. Le budget informatique \niveauexpert}]
Le bilan économique du schéma directeur de LiZbiZ n'a pas été remis
(\NEK{}). Concevez trois tests qui permettraient de lever cette incertitude
sans ce document, et dites ce que chacun établirait.
\end{boxexo}

\begin{boxexo}[title={\ding{45}\quad Exercice~17.10. Une gouvernance de façade \niveauexpert}]
Le comité des risques de la BANQUE SANAGA s'est réuni quatre fois en 2025,
mais n'a jamais été saisi de la bascule du système central. Rédigez
l'observation en sept éléments et sa gravité.
\end{boxexo}

\bigskip

\begin{boxinfo}
\textbf{Points de contrôle de l'élève.} Je sais~: décomposer l'audit de la
fonction informatique~; évaluer un organe par ses traces~; auditer un
schéma directeur et un budget~; repérer une dépendance à une personne~;
recalculer un indicateur et une disponibilité.
\end{boxinfo}

\begin{boxinfo}
\textbf{Pour aller plus loin.}
\begin{itemize}
    \item Le guide d'audit des SI du CHAI~\cite{chai2014guide}~: les fiches
          « pilotage des SI ».
    \item ISACA, COBIT~2019~\cite{isaca2018cobitobjectifs}~: les objectifs
          EDM et APO.
\end{itemize}
\end{boxinfo}

\begin{boxinfo}
\textbf{Transition.} La gouvernance décide~; les applications exécutent.
Le Chapitre~\ref{chap:audit-applications} audite une application en
service~: entrées, traitements, sorties, interfaces, et les tests sur les
données.
\end{boxinfo}
